\structsection{Введение}

Задача упаковки объектов в ограниченный контейнер является одной из классических задач комбинаторной оптимизации. В трехмерном варианте требуется разместить набор прямоугольных параллелепипедов внутри контейнера или на паллете так, чтобы объекты не пересекались и не выходили за заданные границы. В прикладной логистике такая задача возникает при сборке заказов, паллетизации товаров, загрузке транспорта и планировании складских операций.

В крупном ритейле качество паллетизации напрямую влияет на стоимость перевозки и устойчивость товарного потока. Недостаточно плотная укладка приводит к перевозке незаполненного объема, а неучет веса, хрупкости и устойчивости увеличивает риск повреждения товаров. Поэтому практическая постановка 3D Bin Packing существенно шире чисто геометрической задачи: необходимо одновременно учитывать размеры, массу, допустимые ориентации, опору снизу, возможность штабелирования и время расчета.

Тема настоящей выпускной квалификационной работы связана с проектом Smart 3D Packing. Условие проекта формулирует задачу построения алгоритмического модуля 3D Pallet Packing Optimizer, который по описанию паллеты и списка коробов возвращает план размещения с координатами и ориентациями каждого короба. Согласно условию, решение проходит двухэтапную проверку: сначала контролируются жесткие ограничения физической корректности, затем вычисляется взвешенная метрика качества.

Актуальность работы обусловлена тем, что задача 3D Bin Packing относится к NP-трудным задачам~\cite{martello2000}, а практически значимые экземпляры содержат десятки и сотни объектов. Полный перебор быстро становится неприменимым, поэтому в реальных системах используются эвристики, локальный поиск, метаэвристики и гибридные методы~\cite{faroe2003,pisinger2010}. В исследуемом проекте реализовано несколько таких подходов, что позволяет сравнить их в едином экспериментальном контуре.

\textbf{Цель работы} -- исследовать эвристические и ML-подходы к решению задачи 3D Bin Packing в проекте Smart 3D Packing и определить, какие методы обеспечивают лучший баланс качества укладки и времени расчета.

\textbf{Для достижения цели решаются следующие задачи:}
\begin{enumerate}
  \item изучить условие задачи Smart 3D Packing, требования к входным и выходным данным и правила валидации;
  \item формализовать ограничения и метрику качества упаковки;
  \item проанализировать структуру репозитория и восстановить pipeline решения;
  \item описать генерацию синтетических датасетов и состав используемых наборов данных;
  \item рассмотреть реализованные алгоритмы: greedy, MaxRects, LNS, GAN+GA, brute force и многопаллетное расширение;
  \item провести сравнительный анализ сохраненных экспериментальных результатов;
  \item сформулировать ограничения текущей реализации и направления дальнейшего развития.
\end{enumerate}

\textbf{Объект исследования} -- алгоритмическая система трехмерной укладки коробов на паллету.

\textbf{Предмет исследования} -- эвристические, метаэвристические и ML-ориентированные методы построения допустимой и качественной укладки в задаче 3D Bin Packing.

\textbf{Практическая значимость} работы заключается в том, что исследуемая система объединяет генератор данных, набор решателей, валидатор, средства визуализации и сохраненные benchmark-результаты. Такая структура может быть использована как основа для дальнейших исследований алгоритмов паллетизации и для прототипирования прикладного логистического модуля.

Работа состоит из введения, четырех разделов основной части, заключения, списка использованных источников и приложений. В первом разделе рассматриваются предметная область и постановка задачи. Во втором разделе приведен обзор методов решения 3D Bin Packing. В третьем разделе описана архитектура проекта Smart 3D Packing и реализация решателей. В четвертом разделе представлены датасеты, эксперименты, метрики и интерпретация результатов.
